\section{Conclusion}
Using published data and a transparent, reproducible model, LFP batteries outperform lead-acid for Bangladesh easy-bikes on efficiency (\effLFP{} vs \effLA\%), life (about \lifeLFP{} vs \lifeLA{} months), 10-year cost (BDT \tcoLFP{} vs \tcoLA) and carbon (about \saveLFPbike{}~kg\,CO$_2$e per bike-year). NMC lies between the two. The ranking is the same at Khulna, Chuadanga and Rajshahi, so site temperature does not change it, while battery heating and heat waves widen the lead-acid disadvantage. A full switch of a 1.4~M fleet would save about \fleetLFPa{}~Mt\,CO$_2$e per year and remove \leadLow--\leadHigh{}~kt of lead per year from informal recycling. The premium is repaid in about a year, pointing policy toward financing, taxes and recycling rules. LFP's cost advantage is robust to any single input but can disappear when several unfavourable assumptions combine; operator records on distance, pack temperature and scrap value are the priority for confirming it.
\section{Reproducibility}
Run in MATLAB R2024a: (1) \texttt{easybike\_main} (tables and figures); (2) \texttt{easybike\_co2} (carbon, lead and policy); (3) \texttt{easybike\_build\_simulink} (Simulink verification; needs Simulink); (4) \texttt{easybike\_check} (compares your \texttt{results\_summary.csv} with the values in this paper). Inputs live in \texttt{easybike\_params.m}, each tagged with its source key. NASA POWER files: Monthly, T2M, 2024-01-01 to 2025-12-31, local solar time, points 22.84$^\circ$N 89.57$^\circ$E, 23.64$^\circ$N 88.69$^\circ$E and 24.38$^\circ$N 88.60$^\circ$E. Tables and figures in this paper were generated by a Python port that reproduced a previous MATLAB \texttt{results\_summary.csv} to 15 significant digits; the MATLAB code in the final form must be run once to confirm Tables 4--8.
